\documentclass[10pt,a4paper]{article}
\usepackage[utf8]{inputenc}\usepackage[T1]{fontenc}\usepackage[german]{babel}
\usepackage{graphicx,float,booktabs,xcolor,tabularx,enumitem}
\usepackage[margin=2.2cm]{geometry}\usepackage[hidelinks]{hyperref}
\hypersetup{pdfauthor={Thomas Schmiegelt},pdftitle={Was KI-Agenten im Maschinenbau heute schon können}}
\newcommand{\bild}[3]{\begin{figure}[H]\centering\includegraphics[width=#2\linewidth]{bilder/#1}\caption{#3}\end{figure}}
\newcommand{\zwei}[3]{\begin{figure}[H]\centering\includegraphics[width=0.49\linewidth]{bilder/#1}\hfill
  \includegraphics[width=0.49\linewidth]{bilder/#2}\caption{#3}\end{figure}}
\newcommand{\kasten}[2]{\par\smallskip\noindent\fcolorbox{gray}{gray!8}{\parbox{0.96\linewidth}{\textbf{#1} #2}}\par\smallskip}
\newcommand{\zitat}[1]{\par\smallskip\noindent\hspace*{1em}\parbox{0.92\linewidth}{\itshape „#1“}\par\smallskip}
\title{\textbf{Was KI-Agenten im Maschinenbau heute schon können}\\[2pt]\large Ein Selbstversuch:
  vier Tage, ein Rechner, zwei Agenten -- und die Spulen eines Fusionsreaktors}
\author{Thomas Schmiegelt \textbar{} Hobbyprojekt \texttt{stellarator-design}}
\date{25.10.2026}
\begin{document}\maketitle

\kasten{Hobbyprojekt, aus Neugier entstanden.}{Ich bin kein Fusionsforscher. Keine Verbindung zu einem
Forschungsinstitut oder Unternehmen, nichts davon ist begutachtet. Das technische Modell ist an Wendelstein~7-X
\emph{angelehnt} und stark vereinfacht. Dieser Text handelt nicht von den Ergebnissen, sondern von der
Arbeitsweise. Die Zahlen stehen in den beiden technischen Berichten im selben Ordner.}

\section*{Worum es geht}
Kann man als Ingenieur mit KI-Agenten ein Thema bearbeiten, von dem man wenig weiß und für das man eigentlich
ein Team bräuchte -- Physik, CAD, Vernetzung, FEM, Auswertung, Dokumentation? Ich habe es an einem bewusst
schwierigen Bauteil ausprobiert: den dreidimensional verwundenen Magnetspulen eines Stellarators, also eines
Fusionsreaktors, der sein Plasma allein mit äußeren Spulen einschließt \cite{bosch2013}. Nach vier Tagen stand
eine durchgängige Kette von der Spulenbahn bis zur Spannung im Gehäuse. Interessanter als das Ergebnis ist, was
unterwegs alles ging -- und was nicht.

\bild{w7x_torus3d_blender.png}{0.92}{Ergebnis nach vier Tagen: 70 Spulen mit Gehäusen und Magnetfeldlinien,
gerechnet mit freier Software, gerendert in Blender. Gebaut habe ich davon selbst keine Zeile Code.}

%======================================================================
\section{Das Team}
\begin{table}[H]\centering\small
\begin{tabularx}{\linewidth}{@{}lX@{}}\toprule
\textbf{Ich} & Ziel, Richtung, Entscheidungen. Ich schaue mir Bilder an, stelle Fragen und bringe Ideen aus der
  Konstruktionspraxis ein. \\
\textbf{Hermes Agent} \cite{hermes} & Läuft lokal auf meinem Rechner mit offenen Sprachmodellen. Recherche,
  Literatur, Projektordner pflegen, lange Rechenläufe über Nacht. \\
\textbf{Claude Code} \cite{claudecode} & Programmiert: Geometrie, Vernetzung, FEM-Eingabe, Auswertung, Bilder.
  Sucht Fehler. \\
\textbf{Freie Software} & FreeCAD \cite{freecad}, gmsh \cite{gmsh}, ElmerFEM \cite{elmer}, ParaView
  \cite{paraview}, Blender \cite{blender}, Fachbibliotheken aus der Fusionsforschung \cite{simsopt,desc}. \\\bottomrule
\end{tabularx}\end{table}
Die Agenten bedienen diese Programme nicht über die Oberfläche, sondern über Skripte: FreeCAD ohne Fenster,
Elmer über Eingabedateien, Blender im Hintergrund. Ich öffne das Ergebnis dann in der gewohnten Oberfläche und
kann dort weiterarbeiten.

%======================================================================
\section{Was die Agenten konkret gemacht haben}
\begin{table}[H]\centering\small
\begin{tabularx}{\linewidth}{@{}p{3.6cm}X@{}}\toprule
Recherche & Grundlagen erarbeitet; 33 PDFs aus meiner Sammlung gesichtet und die vier relevanten ausgewertet;
  gezielt auf arXiv gesucht, mit Pausen, um nicht gesperrt zu werden. \\
Werkzeuge einrichten & FEM-Löser installiert und geprüft, FreeCAD eingerichtet, Python-Umgebungen für
  unterschiedliche Bibliotheken getrennt. \\
Validierung & Zuerst einen Fall mit bekannter Lösung gerechnet (gerader Leiter), bevor es ans eigentliche
  Bauteil ging. \\
Daten finden & Die echten Spulenbahnen von W7-X in einer frei verfügbaren Forschungsbibliothek gefunden
  \cite{simsopt}. \\
Parametrisches CAD & Ein FreeCAD-Modell erzeugt, das ich in der Oberfläche weiterbearbeiten kann:
  Parametertabelle, Skizzenmaße als Ausdrücke, Symmetrien als Klone, Verbinder. \\
Vernetzung und FEM & Netze erzeugt, Magnetfeld und Spannung in einem gekoppelten Lauf gerechnet, die Kraft mit
  einer zweiten Methode gegengeprüft. \\
Prüfroutinen & Kollisionsprüfung Spule für Spule, mit Video des Aufbaus. \\
Auswertung & Bilder in pyvista, ParaView und Blender; Querschnitte, Schnittmodelle, Detailmodell des Leiters. \\
Dokumentation & Befund-Protokoll, Berichte als PDF, Veröffentlichung auf GitHub. \\\bottomrule
\end{tabularx}\end{table}

\zwei{freecad_gui_w7x.png}{aufbau_schritt35.png}{Links: von einem Agenten per Skript erzeugtes, parametrisches
FreeCAD-Modell, geöffnet in der normalen Oberfläche. Rechts: automatische Aufbauprüfung, Spule für Spule, mit
Kollisionszähler.}

%======================================================================
\section{So sieht die Zusammenarbeit aus}
Der Umgangston ist der eines Gesprächs mit einem Kollegen. Ich schreibe, was ich will, oft knapp und mit
Tippfehlern -- der Agent setzt es um, zeigt ein Bild, und ich entscheide, wie es weitergeht. Einige meiner
tatsächlichen Anweisungen aus dem Projekt (Rechtschreibung korrigiert):
\zitat{Lass das erst sein. Mach erst die Recherche zur Funktion, zum Aufbau und zur Physik.}
\zitat{Du kannst natürlich ein Logiksystem bauen, in dem du die benachbarten Spulen mit einem kleinen Offset
abziehst.}
\zitat{Wir können das auch parallel machen: den funktionierenden Weg für die Berechnung von heute Morgen und den
parallelen CAD-Weg als Add-on.}
\zitat{Supraleitungskabel und Kupferkabel kann man nicht unterscheiden. Durch was fließt der Strom?}
Die zweite Anweisung wurde in wenigen Minuten zu einer Konstruktionslogik in FreeCAD: Jedes Gehäuse zieht die
Hüllen seiner Nachbarn ab. Danach war das Modell kollisionsfrei. Die letzte Frage führte zu einem eigenen
3D-Detailmodell des Leiters, das die Antwort zeigt.
\zwei{w7x_spule_querschnitt.png}{w7x_leiter3d_blender.png}{Aus Fragen werden Bilder: Schichtaufbau einer Spule und
ein Leiterstück, in dem der Strom im türkisfarbenen Supraleiter fließt.}

\kasten{Was dabei auffällt:}{Die Ideen, die das Projekt aus Sackgassen geholt haben, kamen aus der
Konstruktionserfahrung: ein \emph{Variable Section Sweep} wie in Creo, Nachbarn abziehen, CAD und Rechnung
trennen. Die Agenten haben sie schnell umgesetzt, getestet und dokumentiert. Erfahrung wird also nicht
überflüssig -- sie wird zum Hebel.}

%======================================================================
\section{Wo es hakt}
\subsection{Plausibel heißt nicht richtig}
Das erste Spulenmodell sah gut aus und lieferte Zahlen in der richtigen Größenordnung. Es war trotzdem kein
Stellarator: Ein Teil der Spulen umschloss das Plasma gar nicht. Erst eine einfache geometrische Prüfung deckte
es auf. Später fanden sich ein Vorzeichenfehler und eine falsche Stromrichtung, die der Löser ohne Warnung selbst
gewählt hatte.
\zwei{v01_spulen_iso.png}{w7x_wickelpakete.png}{Links das erste, falsche Modell -- rechts die echten Spulen. Den
Unterschied sieht man erst, wenn man weiß, worauf man achten muss.}

\subsection{Quellen sind nicht automatisch korrekt}
Eine vom Agenten erstellte Literaturliste enthielt Einträge ohne Autoren und mit falschen Zeitschriften. Für die
Veröffentlichung wurde jede Quelle einzeln gegen Volltext und Verlagsdaten geprüft.

\subsection{Ein gültiges CAD-Modell ist noch kein Rechenmodell}
Das parametrische FreeCAD-Modell war schön und gültig, aber der Vernetzer scheiterte an Flächen, die mehrere
Boolesche Operationen gemeinsam erzeugt hatten. Die Lösung war ein zweiter Weg, der das Rechenmodell ohne
CAD-Kern direkt aus den Daten erzeugt. Das kennt jeder Berechner: CAD-Modell und FE-Modell sind zwei verschiedene
Dinge.
\bild{freecad_und_elmer.png}{0.92}{Zwei Wege nebeneinander: parametrisches CAD für die Konstruktion (links),
direkt erzeugtes Rechenmodell für die Spannung (rechts).}

\subsection{Gedächtnis und Ressourcen}
Agenten vergessen: Lange Sitzungen werden zusammengefasst, Details gehen verloren. Deshalb steht alles Wichtige in
Dateien -- ein Befund-Protokoll, eine Übergabedatei, kurze Anleitungen („Skills“), die der Agent vor der nächsten
Aufgabe liest. Und Rechenzeit bleibt Rechenzeit: Große FEM-Läufe brauchen Speicher, laufen über Nacht und müssen
so gestartet werden, dass sie ein Sitzungsende überleben.

%======================================================================
\section{Meine Arbeitsregeln}
\begin{enumerate}[itemsep=1pt]
\item \textbf{Erst verstehen, dann bauen.} Recherche vor Werkzeug.
\item \textbf{Ergebnisse als Bild zeigen lassen.} Ein Bild entlarvt Fehler schneller als jede Tabelle.
\item \textbf{Alles zweimal rechnen lassen}, mit unabhängiger Methode. Einfache Prüfungen (Abstände, Umlaufsinn,
  Gleichgewicht der Kräfte) automatisch einbauen.
\item \textbf{Mit einem bekannten Fall anfangen}, bevor das eigentliche Bauteil gerechnet wird.
\item \textbf{Den funktionierenden Weg nicht aufgeben}, während man den eleganten erkundet.
\item \textbf{Befunde und Entscheidungen aufschreiben lassen} -- auch die Fehlschläge mit Ursache.
\item \textbf{Quellen selbst prüfen}, bevor etwas veröffentlicht wird.
\item \textbf{Die Verantwortung bleibt beim Ingenieur.} Kein Ergebnis ohne eigene Plausibilitätsprüfung
  weitergeben.
\end{enumerate}

%======================================================================
\section{Was man damit alles machen kann}
Aus diesem Selbstversuch lassen sich Anwendungen ableiten, die in jedem Konstruktions- oder Berechnungsbüro
denkbar sind:
\begin{itemize}[itemsep=1pt]
\item \textbf{Einarbeitung in fremde Themen:} Grundlagen, Normen- und Literaturlage in Stunden statt Wochen --
  mit Prüfpflicht.
\item \textbf{CAD automatisieren:} parametrische Modelle per Skript erzeugen, Varianten durchspielen, Klone und
  Symmetrien konsequent nutzen; das Ergebnis bleibt im gewohnten CAD-Programm bearbeitbar.
\item \textbf{Berechnungsketten aufbauen:} Geometrie, Netz, Löser und Auswertung verbinden, auch zwischen
  Programmen, die nichts voneinander wissen.
\item \textbf{Variantenstudien über Nacht:} Parameter ändern, rechnen, auswerten, vergleichen -- der Agent startet
  und überwacht, der Ingenieur bewertet am Morgen.
\item \textbf{Prüfroutinen:} Kollisionen, Abstände, Massen, Kräftegleichgewicht automatisch nach jedem Schritt.
\item \textbf{Visualisierung und Kommunikation:} Schnittbilder, Explosions- und Detailmodelle, Renderings für
  Präsentationen, Berichte als PDF.
\item \textbf{Wissen sichern:} Befund-Protokolle und Anleitungen, die auch der nächste Kollege -- oder der nächste
  Agent -- lesen kann.
\end{itemize}

\section*{Fazit}
KI-Agenten ersetzen keine Ingenieurserfahrung. Sie verschieben die Arbeit: weniger Tippen, Suchen und
Formatieren, mehr Fragen stellen, Ergebnisse beurteilen und Entscheidungen treffen. Wer sauber prüft, kann mit
einem einzigen Rechner und freier Software Dinge ausprobieren, für die früher ein Team und Wochen nötig gewesen
wären. Wer nicht prüft, bekommt sehr schnell sehr überzeugende falsche Ergebnisse.

\medskip\noindent Technische Berichte, Programmcode und Befundprotokoll:
\url{https://github.com/ThomasSchmiegelt/stellarator-design}

\begin{thebibliography}{99}\small
\bibitem{bosch2013} H.-S. Bosch, R.\,C. Wolf, T. Andreeva et al.: Technical challenges in the construction of
  the steady-state stellarator Wendelstein 7-X. \emph{Nucl.\ Fusion} 53 (2013) 126001.
  doi:10.1088/0029-5515/53/12/126001
\bibitem{hermes} Nous Research: Hermes Agent. \url{https://github.com/NousResearch/hermes-agent} (MIT-Lizenz).
\bibitem{claudecode} Anthropic: Claude Code. \url{https://www.anthropic.com/claude-code}
\bibitem{simsopt} M. Landreman, B. Medasani, F. Wechsung, A. Giuliani, R. Jorge, C. Zhu: SIMSOPT: A flexible
  framework for stellarator optimization. \emph{J.\ Open Source Softw.} 6(65) (2021) 3525.
  doi:10.21105/joss.03525
\bibitem{desc} D.\,W. Dudt, E. Kolemen: DESC: A stellarator equilibrium solver. \emph{Phys.\ Plasmas} 27 (2020)
  102513.
\bibitem{freecad} FreeCAD, Version 1.1. \url{https://www.freecad.org}
\bibitem{gmsh} C. Geuzaine, J.-F. Remacle: Gmsh: A 3-D finite element mesh generator with built-in pre- and
  post-processing facilities. \emph{Int.\ J.\ Numer.\ Meth.\ Eng.} 79 (2009) 1309--1331.
\bibitem{elmer} ElmerFEM, Version 9.0. CSC -- IT Center for Science, Finnland. \url{https://www.elmerfem.org}
\bibitem{paraview} J. Ahrens, B. Geveci, C. Law: ParaView: An end-user tool for large data visualization. In:
  \emph{Visualization Handbook}, Elsevier (2005).
\bibitem{blender} Blender, Version 4.0. Blender Foundation. \url{https://www.blender.org}
\end{thebibliography}
\end{document}
